% =====================================================================
% Group B -- Chapter 5: orbital motion in multibody environments (one question per exam)
% Source: latex_notes/SFM_05_Orbital Motion in Multibody Environments/cap05.tex
% =====================================================================
\qgroup{B}{Chapter 5 -- Multibody environments and the CR3BP}

One of the three orbital-mechanics questions is always on Chapter 5. Notation of the notes: primaries $m_1>m_2$, mass
parameter $\mu=m_2/(m_1+m_2)$, synodic coordinates $(x,y,z)$, angular velocity $\omega$ of the primaries ($\omega=1$\,TU$^{-1}$
in canonical units, as in the exam texts). Earth--Moon: $\mu=1/82.27=0.01216$.

\begin{center}
\small
\begin{tabular}{@{}lllr@{}}
\toprule
Q & Topic & Asked & Page\\
\midrule
\ref{q:B1} & Jacobi interval, Earth--Moon transfers without escape & $\times$2, Jan 2023, Jul 2025 & \pageref{q:B1}\\
\ref{q:B2} & Synodic frame, 5 libration points, ZVC with/without escape & $\times$2, Feb 2025 & \pageref{q:B2}\\
\ref{q:B3} & Collinear libration points: equation for their position & list & \pageref{q:B3}\\
\ref{q:B4} & Proof of the Jacobi integral & $\times$2, Sep 2024, Jan 2025, Mar 2025 & \pageref{q:B4}\\
\ref{q:B5} & Synodic frame: primaries and their $\mu$-dependence & list & \pageref{q:B5}\\
\ref{q:B6} & $N$-body integrals, Laplace plane, $\underline H_c$ conserved & $\times$2, Jun 2025, Sep 2025 & \pageref{q:B6}\\
\bottomrule
\end{tabular}
\end{center}

Common data used in this group (computed from $\Omega$ with $\mu=1/82.27$, $\omega=1$):
\begin{center}
\small
\begin{tabular}{@{}lccccc@{}}
\toprule
 & $L_1$ & $L_2$ & $L_3$ & $L_{4,5}$\\
\midrule
$x$ [DU] & $0.8369$ & $1.1557$ & $-1.0051$ & $0.5-\mu$, $y=\pm\sqrt3/2$\\
$C_i=2\Omega(L_i)$ & $3.18838$ & $3.17220$ & $3.01215$ & $2.98799$\\
\bottomrule
\end{tabular}
\end{center}
{\footnotesize ($C_1$ in the notes reads 3.18338: typo corrected, errata \#20.) DU $=R=384400$\,km, TU $=1/\omega\simeq4.34$\,days.}

% ---------------------------------------------------------------------
\question{B1}{Earth--Moon CR3BP: Jacobi interval for transfers without escape}
% source: cap05.tex l.379-467 (Jacobi integral, ZVS), l.469-571 (libration points, Omega at Li), l.573-613 (special values of C)
\begin{qbox}
In the circular restricted three-body problem (CR3BP) with Earth and Moon as primaries, identify the Jacobi-integral
interval that allows Earth--Moon transfers without any possible escape toward the outer space, referencing the
zero-velocity curves. Depict the curve geometry.
\qmeta{$\times$2 in the list (2025); Jan 2023; Jul 2025 ("provide the definition of the Jacobi integral").}
\qnote{Difference from \qref{B2}: B2 also asks for the synodic frame, the five libration points and the case \emph{with} an
escape route.}
\end{qbox}

\begin{outline}
\begin{enumerate}
\item CR3BP hypotheses; synodic frame; potential $\Omega$; equations $\ddot x-2\omega\dot y=\Omega_x$, \dots
\item Jacobi integral $C=2\Omega-(\dot x^2+\dot y^2+\dot z^2)$ (one-line proof); $C$ decreases as energy increases.
\item Zero-velocity surfaces/curves: $v^2=2\Omega-C\ge0$ defines the Hill region; shape (spheres near primaries, cylinder far away).
\item Libration points and $C_i=2\Omega(L_i)$: $C_1>C_2>C_3>C_{4,5}$; Earth--Moon values.
\item Answer: $C_2<C<C_1$: the neck at $L_1$ is open (Earth--Moon transfer), the neck at $L_2$ closed (no escape). Sketch.
\end{enumerate}
\end{outline}

\answer
\paragraph{Model.} The CR3BP studies a third body (spacecraft) of negligible mass, $m\ll m_2<m_1$, attracted by two
\emph{primaries} (Earth and Moon) that move on circular orbits about their centre of mass with angular velocity
$\omega=\sqrt{G(m_1+m_2)/R^3}$; the spacecraft does not affect the primaries. In the \textbf{synodic frame} (origin at the centre of
mass, $\hat\imath$ from $m_1$ towards $m_2$, $\hat k$ along the angular momentum of the primaries, rotating with $\omega$) the primaries are
fixed at $x_1=-\mu$, $x_2=1-\mu$ (DU $=R$, TU $=1/\omega$), and the motion obeys
\begin{equation}
\ddot{x}-2\omega\dot{y}=\frac{\partial\Omega}{\partial x},\qquad
\ddot{y}+2\omega\dot{x}=\frac{\partial\Omega}{\partial y},\qquad
\ddot{z}=\frac{\partial\Omega}{\partial z},
\end{equation}
\begin{equation}
\Omega=\frac{\omega^{2}}{2}(x^{2}+y^{2})+\frac{1-\mu}{r_1}+\frac{\mu}{r_2},\qquad
r_1=\sqrt{(x+\mu)^{2}+y^{2}+z^{2}},\ \ r_2=\sqrt{(x+\mu-1)^{2}+y^{2}+z^{2}} .
\end{equation}

\paragraph{Jacobi integral (definition).} Multiplying the three equations by $\dot x$, $\dot y$, $\dot z$ and adding, the Coriolis terms
cancel and, since $\Omega$ depends only on the position, $\frac12\frac{d}{dt}(\dot x^2+\dot y^2+\dot z^2)=\frac{d\Omega}{dt}$. Hence
\begin{equation}
C:=2\Omega(x,y,z)-\left(\dot{x}^{2}+\dot{y}^{2}+\dot{z}^{2}\right)=\text{const}
\end{equation}
is an integral of motion of the CR3BP, fixed by the initial conditions. It is related to the energy: \textbf{$C$ decreases as the
mechanical energy increases} (a $\Delta v$ that increases the speed lowers $C$).

\paragraph{Zero-velocity surfaces and curves.} Since $v^2=2\Omega(x,y,z)-C\ge0$, motion is possible only where $2\Omega\ge C$
(\emph{Hill's region}). The boundary $2\Omega=C$, where the velocity would be zero, is the zero-velocity surface (curve in the
$x$--$y$ plane). For large $x,y$ the term $\omega^2(x^2+y^2)$ dominates (near-cylinder about $z$); close to a primary the
corresponding $1/r_i$ term dominates (near-sphere around $m_1$ or $m_2$). For large $C$ (low energy) motion is confined
inside the spheres or outside the cylinder; as $C$ decreases the regions grow and \textbf{merge through "necks" located at the
libration points}.

\paragraph{Critical values of $C$.} At a libration point $L_i$ (equilibrium in the synodic frame) the velocity is zero, so $L_i$
lies on the zero-velocity curve with $C_i=2\Omega(L_i)$. For Earth--Moon:
$C_1=3.18838>C_2=3.17220>C_3=3.01215>C_{4,5}=2.98799$ (canonical units). The topology changes exactly at these values:
\begin{itemize}
\item $C>C_1$: three separate regions (around the Earth, around the Moon, outside): no transfer possible;
\item $C_2<C<C_1$: the neck at $L_1$ opens;
\item $C_3<C<C_2$: also the neck at $L_2$ opens (escape);
\item $C_{4,5}<C<C_3$: forbidden only near $L_4$, $L_5$; \quad $C<C_{4,5}$: motion allowed everywhere.
\end{itemize}

\paragraph{Answer: Earth--Moon transfers without escape.}
\begin{equation}
\boxed{\;C_2<C<C_1\qquad(3.17220<C<3.18838\ \text{for Earth--Moon})\;}
\end{equation}
The inner zero-velocity curve has an opening at $L_1$, between Earth and Moon: an \emph{interior transfer} from the region around
the Earth to the region around the Moon is feasible, passing close to $L_1$. The opening at $L_2$ (behind the Moon) is still closed:
the inner region (Earth + Moon) is separated by a forbidden ring from the exterior region, so \textbf{no escape} toward outer space is
possible. Motion outside the outer curve is also allowed, but it cannot be reached from the inner region.

\begin{center}
\includegraphics[width=0.48\linewidth]{zvc_noescape}\hfill
\begin{minipage}[b]{0.48\linewidth}\small
Computed with $\mu=1/82.27$ and $C=(C_1+C_2)/2$. Hatched: forbidden region ($2\Omega<C$).
The white region around E and M is connected through the $L_1$ neck;
the hatched ring is closed at $L_2$, so the exterior (white, outside) is unreachable.\\[4pt]
Lowering $C$ (more energy, e.g.\ a larger injection $\Delta v$) widens the $L_1$ neck; at $C=C_2$ the $L_2$ neck opens and escape
becomes possible.
\end{minipage}
\end{center}

\begin{fig2draw}
E (big dot) and M (small dot) on the $x$-axis. Draw a closed curve around the Earth and a small one around the Moon touching
it at $L_1$ (between them, slightly open: a "peanut" with a neck), and a large closed oval outside both, closed near $L_2$.
Hatch the ring between the inner peanut and the outer oval. Write $C_2<C<C_1$.
\end{fig2draw}

\begin{keyf}
$C=2\Omega-(\dot x^2+\dot y^2+\dot z^2)$, \ $\Omega=\frac{\omega^2}{2}(x^2+y^2)+\frac{1-\mu}{r_1}+\frac{\mu}{r_2}$, \
Hill region $2\Omega\ge C$, \ $C_i=2\Omega(L_i)$, \ $C_1>C_2>C_3>C_{4,5}$, \ no-escape transfer: $C_2<C<C_1$.
\end{keyf}

\pitfalls
\begin{itemize}
\item Higher $C$ means \emph{lower} energy: the order of the intervals is reversed with respect to the energy.
\item "Without escape" requires the $L_2$ neck closed (lower bound $C_2$) and "transfer" requires the $L_1$ neck open (upper bound $C_1$).
\item The forbidden region is where $2\Omega<C$; hatch it clearly and mark $L_1$, $L_2$.
\end{itemize}

\source{cap05.tex l.244--305 (CR3BP, synodic frame), l.307--413 (equations, Jacobi integral), l.415--467 (zero-velocity surfaces),
l.573--613 (special values of $C$, Earth--Moon values). Figure computed in \texttt{figures/src/make\_figures.py}.}

% ---------------------------------------------------------------------
\question{B2}{Synodic frame, the five libration points, zero-velocity curves with and without escape}
% source: cap05.tex l.259-305, 469-535, 573-613
\begin{qbox}
In the CR3BP, plot the five libration points in the synodic frame and define that frame. Then, for Earth--Moon primaries,
draw the zero-velocity curves that allow Earth-to-Moon transfers without escape; also show the case where an escape route
exists, and state the corresponding Jacobi-constant interval.
\qmeta{$\times$2 in the list (2025); Feb 2025 ("draw the case with a single escape route outward and define the interval of $C$").}
\qnote{Difference from \qref{B1}: adds the synodic frame, the libration points and the escape case.}
\end{qbox}

\begin{outline}
\begin{enumerate}
\item Synodic frame: origin at the centre of mass, $\hat\imath$ towards $m_2$, $\hat k\parallel\underline H$, rotating with $\omega$; canonical units; primaries at $-\mu$, $1-\mu$.
\item Libration points: equilibria with zero velocity, $\Omega_x=\Omega_y=0$; three collinear ($L_3<x_1<L_1<x_2<L_2$) and two triangular ($r_1=r_2=1$). Sketch.
\item Jacobi integral and Hill region $2\Omega\ge C$; $C_i=2\Omega(L_i)$, $C_1>C_2>C_3>C_{4,5}$.
\item No escape: $C_2<C<C_1$ (sketch). Single escape route: $C_3<C<C_2$, exit through $L_2$ (sketch).
\item Complete list of the five regimes.
\end{enumerate}
\end{outline}

\answer
\paragraph{Synodic frame.} Two primaries $m_1>m_2$ move on circular orbits about their centre of mass $O'$ with constant
distance $R$ and angular velocity $\omega=\sqrt{G(m_1+m_2)/R^3}$. The synodic frame $\{\hat\imath,\hat\jmath,\hat k\}$ has origin in
$O'$, $\hat\imath$ along the line of the primaries (from $m_1$ to $m_2$), $\hat k$ parallel to their angular momentum
$\underline H$, $\hat\jmath=\hat k\times\hat\imath$, and \textbf{rotates with the primaries}: with respect to the inertial axes
$\{\hat c_1,\hat c_2,\hat c_3\}$ (coincident at $t=0$)
\begin{equation}
\begin{bmatrix}\hat\imath\\ \hat\jmath\\ \hat k\end{bmatrix}=
\begin{bmatrix}\cos\omega t&\sin\omega t&0\\-\sin\omega t&\cos\omega t&0\\0&0&1\end{bmatrix}
\begin{bmatrix}\hat c_1\\ \hat c_2\\ \hat c_3\end{bmatrix}.
\end{equation}
In canonical units (DU $=R$, TU $=1/\omega$, so $G(m_1+m_2)=1$) the primaries are \emph{fixed} at $x_1=-\mu$, $x_2=1-\mu$,
with $\mu=m_2/(m_1+m_2)<1/2$. The spacecraft position is $\underline r=x\hat\imath+y\hat\jmath+z\hat k$ and its equations of motion
are $\ddot x-2\omega\dot y=\Omega_x$, $\ddot y+2\omega\dot x=\Omega_y$, $\ddot z=\Omega_z$ with
$\Omega=\frac{\omega^2}{2}(x^2+y^2)+\frac{1-\mu}{r_1}+\frac{\mu}{r_2}$.

\paragraph{Libration (Lagrange) points.} Equilibrium points in the synodic frame, where the spacecraft can stay indefinitely if
placed with zero velocity. They lie in the $x$--$y$ plane and require $\dot x=\dot y=\ddot x=\ddot y=0$, i.e.\ $\Omega_x=\Omega_y=0$:
\begin{itemize}
\item \textbf{collinear points} ($y=0$): one root of a quintic in each interval: $L_3$ ($x<-\mu$, beyond $m_1$), $L_1$ (between the
primaries), $L_2$ (beyond $m_2$) -- see \qref{B3};
\item \textbf{triangular points} ($y\neq0$): $r_1=r_2=1$ DU, i.e.\ $L_4$, $L_5$ form \textbf{equilateral triangles} with the primaries
($x=\frac12-\mu$, $y=\pm\frac{\sqrt3}{2}$), above and below the $x$-axis, for any $\mu$.
\end{itemize}
\begin{center}
\includegraphics[width=0.5\linewidth]{lagrange}\hfill
\includegraphics[width=0.36\linewidth]{synodic}
\end{center}
$\Omega$ is stationary at every $L_i$: only stationary (saddle) at $L_{1,2,3}$, a minimum at $L_{4,5}$ with
$2\Omega_{\min}=3-\mu(1-\mu)$.

\paragraph{Zero-velocity curves.} The Jacobi integral $C=2\Omega-v^2$ is constant; motion is allowed only where $2\Omega\ge C$.
At $L_i$ the velocity is zero, so the curve through $L_i$ corresponds to $C_i=2\Omega(L_i)$, with
$C_1>C_2>C_3>C_{4,5}$ (Earth--Moon: 3.18838, 3.17220, 3.01215, 2.98799).

\paragraph{The two requested cases.}
\begin{center}
\includegraphics[width=0.46\linewidth]{zvc_noescape}\hfill
\includegraphics[width=0.46\linewidth]{zvc_escape}
\end{center}
\begin{itemize}
\item \textbf{Earth--Moon transfer, no escape}: $\boxed{C_2<C<C_1}$. The neck at $L_1$ is open (interior transfer Earth $\leftrightarrow$
Moon close to $L_1$), the neck at $L_2$ is closed: the forbidden ring separates the Earth--Moon region from outer space.
\item \textbf{Single escape route}: $\boxed{C_3<C<C_2}$. Both necks $L_1$ and $L_2$ are open: interior transfers (through $L_1$) and
exterior transfers (through $L_2$) are feasible, and a spacecraft can leave the Earth--Moon system \emph{only} through the opening
behind the Moon at $L_2$; the forbidden region becomes a horseshoe, still closed at $L_3$.
\end{itemize}

\paragraph{All regimes} (as $C$ decreases, i.e.\ energy increases):
\begin{enumerate}[label=(\arabic*)]
\item $C<C_{4,5}$: motion allowed in the entire space (no zero-velocity curve);
\item $C_{4,5}<C<C_3$: motion forbidden only near $L_4$ and $L_5$;
\item $C_3<C<C_2$: interior (via $L_1$) and exterior (via $L_2$) transfers;
\item $C_2<C<C_1$: interior transfers only; outside motion feasible but not reachable;
\item $C>C_1$: no transfer: confined near $m_1$, near $m_2$, or outside.
\end{enumerate}
\begin{center}
\includegraphics[width=0.95\linewidth]{zvc_cases}
\end{center}

\begin{fig2draw}
Axes $x,y$; E at $-\mu$ (left of $O'$, almost on it), M at $1-\mu$; $L_1$ between them close to the Moon, $L_2$ just beyond the Moon,
$L_3$ at $x\approx-1$, $L_4$/$L_5$ at the vertices of the equilateral triangles. Then two sketches: hatched ring closed at
$L_2$ ($C_2<C<C_1$) and hatched horseshoe open at $L_2$ ($C_3<C<C_2$).
\end{fig2draw}

\begin{keyf}
Synodic: origin c.m., $\hat\imath$ $m_1\!\to\! m_2$, $\hat k\parallel\underline H$, rotation $\omega t$; $x_1=-\mu$, $x_2=1-\mu$; \
$L_{4,5}$: $r_1=r_2=1$; \ $C=2\Omega-v^2$; \ no escape $C_2<C<C_1$; escape through $L_2$: $C_3<C<C_2$.
\end{keyf}

\pitfalls
\begin{itemize}
\item $L_1$ is between the primaries, $L_2$ beyond the smaller one, $L_3$ beyond the larger one (on the opposite side).
\item The frame rotates: libration points are equilibria \emph{in the synodic frame} (in inertial space they rotate with the primaries).
\item "Single escape route" = only $L_2$ open (not $L_3$): lower bound $C_3$.
\end{itemize}

\source{cap05.tex l.259--305 (synodic frame), l.469--535 (libration points), l.537--571 ($\Omega$ at $L_i$), l.573--613 (special
values of $C$; errata \#20). Figures computed in \texttt{figures/src/make\_figures.py}.}

% ---------------------------------------------------------------------
\question{B3}{Collinear libration points: conditions and equation for their position}
% source: cap05.tex l.469-508 (+ l.510-535 triangular)
\begin{qbox}
In the CR3BP, draw the libration points in the synodic frame. Starting from the equations of motion
$\ddot x-2\dot y=\Omega_x$, $\ddot y+2\dot x=\Omega_y$, $\ddot z=\Omega_z$, $\Omega=\frac{x^2+y^2}{2}+\frac{1-\mu}{r_1}+\frac{\mu}{r_2}$,
define the conditions of the collinear libration points and derive the equation for their position ($r_1$ and $r_2$ are the
distances to the primaries).
\qmeta{list by chapter (Ch.\,5).}
\end{qbox}

\begin{outline}
\begin{enumerate}
\item Sketch of $L_1\dots L_5$ in the synodic frame.
\item Equilibrium: $\dot x=\dot y=\dot z=0$ and $\ddot x=\ddot y=\ddot z=0$ $\Rightarrow$ $\Omega_x=\Omega_y=\Omega_z=0$; $z=0$.
\item Explicit $\Omega_x$, $\Omega_y$; collinear: $y=0$ $\Rightarrow$ $\Omega_y=0$ automatically.
\item Equation $x-\frac{(1-\mu)(x+\mu)}{r_1^3}-\frac{\mu(x+\mu-1)}{r_2^3}=0$ with $r_1=|x+\mu|$, $r_2=|x+\mu-1|$.
\item Three intervals $\Rightarrow$ three quintics, one real root each: $L_3$, $L_1$, $L_2$; Earth--Moon values.
\end{enumerate}
\end{outline}

\answer
\paragraph{Libration points.} With canonical units ($\omega=1$), primaries at $x_1=-\mu$ (mass $1-\mu$) and $x_2=1-\mu$ (mass $\mu$).
Libration points are equilibrium points in the synodic frame: a spacecraft placed there with zero velocity remains there.
\begin{center}\includegraphics[width=0.5\linewidth]{lagrange}\end{center}

\paragraph{Equilibrium conditions.} Equilibrium requires zero velocity and zero acceleration in the synodic frame:
\begin{equation}
\dot x=\dot y=\dot z=0,\qquad \ddot x=\ddot y=\ddot z=0\;\Longrightarrow\;
\frac{\partial\Omega}{\partial x}=\frac{\partial\Omega}{\partial y}=\frac{\partial\Omega}{\partial z}=0,
\end{equation}
i.e.\ libration points are the stationary points of $\Omega$ (the Coriolis terms vanish because the velocity is zero).
$\Omega_z=-z\left[\frac{1-\mu}{r_1^3}+\frac{\mu}{r_2^3}\right]=0$ gives $z=0$: the points lie in the plane of the primaries. With
$r_1=\sqrt{(x+\mu)^2+y^2}$, $r_2=\sqrt{(x+\mu-1)^2+y^2}$:
\begin{equation}
\sys{\text{(A)}\quad&\Omega_x=x-\frac{(1-\mu)(x+\mu)}{r_1^{3}}-\frac{\mu(x+\mu-1)}{r_2^{3}}=0\\
\text{(B)}\quad&\Omega_y=y-\frac{(1-\mu)\,y}{r_1^{3}}-\frac{\mu\,y}{r_2^{3}}=0}
\end{equation}

\paragraph{Collinear points.} They are sought on the $x$-axis, $y=0$: (B) is then satisfied identically, and only (A) remains,
with $r_1=|x+\mu|$, $r_2=|x+\mu-1|$:
\begin{equation}
\boxed{\;x-\frac{(1-\mu)(x+\mu)}{|x+\mu|^{3}}-\frac{\mu\,(x+\mu-1)}{|x+\mu-1|^{3}}=0\;}
\end{equation}
Physically: gravity of the two primaries balances the centrifugal term $x$ (the point co-rotates with the primaries). The
absolute values split the axis into three intervals; in each one the equation becomes a quintic polynomial in $x$ with a single
admissible real root:
\begin{itemize}
\item[(a)] $x<-\mu$ (beyond $m_1$): $r_1=-(x+\mu)$, $r_2=1-\mu-x$ $\Rightarrow$ \textbf{$L_3$}, left exterior point;
\item[(b)] $-\mu<x<1-\mu$ (between the primaries): $r_1=x+\mu$, $r_2=1-\mu-x$ $\Rightarrow$ \textbf{$L_1$}, interior point;
\item[(c)] $x>1-\mu$ (beyond $m_2$): $r_1=x+\mu$, $r_2=x+\mu-1$ $\Rightarrow$ \textbf{$L_2$}, right exterior point.
\end{itemize}
For example, for $L_1$ with $\gamma=r_2$ (distance from $m_2$), $x=1-\mu-\gamma$, $r_1=1-\gamma$:
\begin{equation}
1-\mu-\gamma-\frac{1-\mu}{(1-\gamma)^2}+\frac{\mu}{\gamma^2}=0,
\end{equation}
which, multiplied by $\gamma^2(1-\gamma)^2$, is a quintic in $\gamma$ \notinnotes. For small $\mu$, $\gamma\approx(\mu/3)^{1/3}$ for both
$L_1$ and $L_2$ (Hill sphere) \notinnotes. Earth--Moon ($\mu=0.01216$): $x_{L_1}=0.8369$, $x_{L_2}=1.1557$, $x_{L_3}=-1.0051$ DU,
i.e.\ $L_1$ about 58\,000\,km and $L_2$ about 64\,500\,km from the Moon.

\paragraph{For completeness: triangular points.} If $y\neq0$, (B) divided by $y$ gives $1-\frac{1-\mu}{r_1^3}-\frac{\mu}{r_2^3}=0$, satisfied by
$r_1=r_2=1$; inserted in (A): $x-(1-\mu)(x+\mu)-\mu(x+\mu-1)=0$ identically. Hence $L_4$, $L_5$ at the vertices of the two
equilateral triangles built on the primaries.

\begin{keyf}
Equilibrium $\Leftrightarrow\nabla\Omega=\underline 0$ ($z=0$); collinear ($y=0$): $x-\dfrac{(1-\mu)(x+\mu)}{r_1^3}-\dfrac{\mu(x+\mu-1)}{r_2^3}=0$,
$r_1=|x+\mu|$, $r_2=|x+\mu-1|$; intervals $x<-\mu$ ($L_3$), $-\mu<x<1-\mu$ ($L_1$), $x>1-\mu$ ($L_2$).
\end{keyf}

\pitfalls
\begin{itemize}
\item The condition is on $\nabla\Omega$, not on the gravitational force alone: the centrifugal term $x$ is essential.
\item Keep the absolute values (or the signs of $x+\mu$ and $x+\mu-1$) in each interval.
\item Naming: $L_1$ between the primaries, $L_2$ beyond $m_2$, $L_3$ beyond $m_1$.
\end{itemize}

\further
Stability (commented section of the notes): $\Omega$ is a saddle at $L_{1,2,3}$, which are \emph{unstable}; $L_{4,5}$ are minima of
$\Omega$ and, according to the linear analysis, neutrally stable if $\mu<\frac12-\frac{\sqrt{69}}{18}\simeq0.0385$ (Earth--Moon
satisfies it).

\source{cap05.tex l.469--508 (collinear points), l.510--535 (triangular points), l.537--571 ($\Omega$ at $L_i$); numerical roots
from \texttt{figures/src/make\_figures.py}.}

% ---------------------------------------------------------------------
\question{B4}{Synodic frame and proof of the Jacobi integral}
% source: cap05.tex l.259-413
\begin{qbox}
In the CR3BP, define the synodic frame and its coordinates $(x,y,z)$. Starting from
$\ddot x-2\dot y=\frac{\partial\Omega}{\partial x}$, $\ddot y+2\dot x=\frac{\partial\Omega}{\partial y}$, $\ddot z=\frac{\partial\Omega}{\partial z}$,
$\Omega=\frac{x^2+y^2}{2}+\frac{1-\mu}{r_1}+\frac{\mu}{r_2}$, prove (and give) the Jacobi integral.
\qmeta{$\times$2 in the list (2025); Sep 2024 ("Jacobi integral"); Jan 2025 ("from the equations with the derivatives of $\Omega$ to the
Jacobi integral"); Mar 2025 ("Jacobi integral").}
\qnote{Sep 2024 Q3 mentions "the same question as January about $\Omega$": if $\Omega$ is the CR3BP potential it is this question;
if it is the RAAN, see \qref{C3}.}
\end{qbox}

\begin{outline}
\begin{enumerate}
\item CR3BP hypotheses; synodic frame (origin, axes, rotation), canonical units, $\mu$, positions of the primaries; coordinates $(x,y,z)$.
\item (Optional) origin of the equations: acceleration in the rotating frame (Coriolis $2\omega$, centrifugal $\omega^2$).
\item Multiply by $\dot x,\dot y,\dot z$ and add: Coriolis terms cancel.
\item $\Omega=\Omega(x,y,z)$ only $\Rightarrow$ RHS $=d\Omega/dt$; LHS $=\frac12 dv^2/dt$ $\Rightarrow$ $C=2\Omega-v^2$ constant.
\item Meaning: integral of motion, related to energy ($C\downarrow$ when $\mathcal E\uparrow$); zero-velocity surfaces.
\end{enumerate}
\end{outline}

\answer
\paragraph{Problem and synodic frame.} CR3BP: a spacecraft of negligible mass moves under the attraction of two primaries
$m_1>m_2$ that describe circular orbits about their centre of mass with constant distance $R$ and angular velocity
$\omega=\sqrt{G(m_1+m_2)/R^3}$. The \textbf{synodic frame} $\{\hat\imath,\hat\jmath,\hat k\}$ has
\begin{itemize}
\item origin $O'$ at the centre of mass of the primaries;
\item $\hat\imath$ along the line joining the primaries, pointing from $m_1$ to $m_2$; $\hat k$ parallel to the angular momentum of the
primaries (normal to their orbit plane); $\hat\jmath=\hat k\times\hat\imath$;
\item it \textbf{rotates} with angular velocity $\underline\omega=\omega\hat k$ with respect to the inertial frame, so the primaries are fixed.
\end{itemize}
Canonical units: DU $=R$, TU $=1/\omega$ (so $\omega=1$ and $G(m_1+m_2)=1$); $\mu=m_2/(m_1+m_2)$, $m_1$ at $x=-\mu$, $m_2$ at
$x=1-\mu$. The coordinates $(x,y,z)$ are the components of $\underline r=x\hat\imath+y\hat\jmath+z\hat k$ in this rotating frame, and
$r_1=\sqrt{(x+\mu)^2+y^2+z^2}$, $r_2=\sqrt{(x+\mu-1)^2+y^2+z^2}$ are the distances from the primaries.

\paragraph{Where the equations come from.} Differentiating $\underline r=x\hat\imath+y\hat\jmath+z\hat k$ twice with
$\dot{\hat\imath}=\omega\hat\jmath$, $\dot{\hat\jmath}=-\omega\hat\imath$:
$\ddot{\underline r}=(\ddot x-2\omega\dot y-\omega^2x)\hat\imath+(\ddot y+2\omega\dot x-\omega^2y)\hat\jmath+\ddot z\hat k$
(Coriolis and centrifugal terms). Equating to the gravitational acceleration of the two primaries and moving the centrifugal terms
to the right-hand side gives the given equations with the pseudo-potential $\Omega$ (centrifugal + gravitational).

\paragraph{Proof.} Multiply the three equations by $\dot x$, $\dot y$, $\dot z$ respectively and add:
\begin{equation}
\dot x\ddot x-2\dot x\dot y+\dot y\ddot y+2\dot y\dot x+\dot z\ddot z
=\dot x\frac{\partial\Omega}{\partial x}+\dot y\frac{\partial\Omega}{\partial y}+\dot z\frac{\partial\Omega}{\partial z}.
\end{equation}
\begin{itemize}
\item The Coriolis terms cancel: $-2\dot x\dot y+2\dot y\dot x=0$ (the Coriolis acceleration is orthogonal to the relative velocity
and does no work).
\item Left-hand side: $\dot x\ddot x+\dot y\ddot y+\dot z\ddot z=\frac12\frac{d}{dt}\left(\dot x^2+\dot y^2+\dot z^2\right)$.
\item Right-hand side: $\Omega$ depends only on $(x,y,z)$, \emph{not explicitly on time} (the primaries are fixed in the synodic frame),
so the chain rule gives $\Omega_x\dot x+\Omega_y\dot y+\Omega_z\dot z=\frac{d\Omega}{dt}$.
\end{itemize}
Therefore
\begin{equation}
\frac12\frac{d}{dt}\left(\dot{x}^{2}+\dot{y}^{2}+\dot{z}^{2}\right)=\frac{d\Omega}{dt}
\;\Longrightarrow\;
\frac{d}{dt}\left[\Omega-\tfrac12\left(\dot{x}^{2}+\dot{y}^{2}+\dot{z}^{2}\right)\right]=0,
\end{equation}
\begin{equation}
\boxed{\;C:=2\Omega(x,y,z)-\left(\dot{x}^{2}+\dot{y}^{2}+\dot{z}^{2}\right)=\text{const}\;}\qquad
2\Omega=x^2+y^2+\frac{2(1-\mu)}{r_1}+\frac{2\mu}{r_2}.
\end{equation}

\paragraph{Meaning.} $C$ (Jacobi constant) is the only known integral of motion of the CR3BP; it is fixed by the initial conditions,
$C=2\Omega(\underline r_0)-v_0^2$ with $v_0$ the velocity \emph{relative to the synodic frame}. It plays the role of an energy:
$C$ \textbf{decreases as the mechanical energy increases} (e.g.\ after a $\Delta v$ that increases the speed). It is not the
inertial energy, which is not conserved because the primaries move. Since $v^2=2\Omega-C\ge0$, it defines the
\textbf{zero-velocity surfaces} $2\Omega=C$, which bound the regions of possible motion (Hill's regions); see \qref{B1}.

\begin{keyf}
$\ddot x-2\dot y=\Omega_x$, $\ddot y+2\dot x=\Omega_y$, $\ddot z=\Omega_z$; \ $\times(\dot x,\dot y,\dot z)$ and sum $\Rightarrow$
$\frac12\frac{d}{dt}v^2=\frac{d\Omega}{dt}$; \ $C=2\Omega-(\dot x^2+\dot y^2+\dot z^2)$, $\Omega=\frac{x^2+y^2}{2}+\frac{1-\mu}{r_1}+\frac{\mu}{r_2}$.
\end{keyf}

\pitfalls
\begin{itemize}
\item The key step is that $\Omega$ is time-independent in the synodic frame; say it explicitly.
\item Show the cancellation of the Coriolis terms.
\item $v$ is the velocity relative to the rotating frame; $C$ is not the energy (it decreases when the energy increases).
\item With dimensional units the equations contain $\omega$ ($\ddot x-2\omega\dot y$, $\frac{\omega^2}{2}(x^2+y^2)$); the proof is identical.
\end{itemize}

\further
In terms of inertial quantities, $C=-2(\mathcal E-\omega H_z)$, where $\mathcal E$ and $H_z$ are the inertial specific energy and angular
momentum component of the spacecraft: neither is conserved separately, but this combination is \notinnotes.

\source{cap05.tex l.244--305 (CR3BP and synodic frame), l.307--378 (equations of motion), l.379--413 (Jacobi integral),
l.415--467 (zero-velocity surfaces). Formulario: not included (theory only).}

% ---------------------------------------------------------------------
\question{B5}{Synodic frame: origin, positions and gravitational parameters of the primaries}
% source: cap05.tex l.111-236 (two-body problem, omega), l.244-305 (synodic frame)
\begin{qbox}
In the CR3BP, define the synodic frame and locate its origin. Show analytically how to obtain the positions of the two primaries in
terms of the mass ratio $\mu$ and give their gravitational parameters in terms of $\mu$.
\qmeta{list by chapter (Ch.\,5).}
\end{qbox}

\begin{outline}
\begin{enumerate}
\item Hypotheses of the CR3BP; the primaries move on circles about their centre of mass $O'$ (two-body problem).
\item Synodic frame: origin $O'$, $\hat\imath$ towards $m_2$, $\hat k\parallel\underline H$, rotating with $\omega=\sqrt{GM/R^3}$; rotation matrix.
\item Positions: $x_2-x_1=R$ and $m_1x_1+m_2x_2=0$ $\Rightarrow$ $x_1=-\frac{m_2}{M}R$, $x_2=\frac{m_1}{M}R$.
\item Canonical units DU $=R$, TU $=1/\omega$ $\Rightarrow$ $GM=$ DU$^3$/TU$^2$; $\mu=m_2/M$: $x_1=-\mu$, $x_2=1-\mu$; $Gm_1=1-\mu$, $Gm_2=\mu$.
\end{enumerate}
\end{outline}

\answer
\paragraph{Motion of the primaries.} In the two-body problem the centre of mass $O'$ of $m_1$ and $m_2$ moves uniformly and can be
taken as origin of an inertial frame; then $m_1\underline R_1+m_2\underline R_2=\underline 0$ and each primary moves on a conic about
$O'$, the two conics having the same eccentricity, with $a_1/a_2=m_2/m_1$ and mean motion $\omega^2=G M/(a_1+a_2)^3$,
$M=m_1+m_2$. In the CR3BP the primaries move on \textbf{circles} ($e=0$) at constant distance $R=a_1+a_2$, hence
\begin{equation}
\omega=\sqrt{\frac{G(m_1+m_2)}{R^{3}}}=\text{const}.
\end{equation}

\paragraph{Synodic frame.} Origin at $O'$ (centre of mass of the primaries); $\hat\imath$ along the line of the primaries pointing towards
$m_2$; $\hat k$ parallel to the angular momentum $\underline H$ of the primaries; $\hat\jmath$ completes the right-handed triad. It
rotates with them: if it coincides with the inertial frame $\{\hat c_i\}$ at $t=0$, $[\hat\imath\ \hat\jmath\ \hat k]^T=R_3(\omega t)[\hat c_1\ \hat c_2\ \hat c_3]^T$.
In this frame the primaries are at rest on the $x$-axis.

\paragraph{Positions of the primaries.} With $x_1<0<x_2$ the abscissas of $m_1$ and $m_2$:
\begin{equation}
\sys{&x_2-x_1=R\\ &m_1x_1+m_2x_2=0\ \text{(origin at the centre of mass)}}
\;\Longrightarrow\;
\sys{&x_1=-\frac{m_2}{m_1+m_2}\,R\\ &x_2=\frac{m_1}{m_1+m_2}\,R}
\end{equation}
(from the second equation $x_2=-\frac{m_1}{m_2}x_1$; substituting in the first: $-x_1\frac{m_1+m_2}{m_2}=R$).

\paragraph{Mass parameter and canonical units.} Define
\begin{equation}
\mu:=\frac{m_2}{m_1+m_2}\in\left(0,\tfrac12\right)\quad(m_1>m_2),\qquad \text{DU}=R,\qquad \text{TU}=\omega^{-1}.
\end{equation}
Then $\omega=1$ TU$^{-1}$ and, from $\omega^2=G(m_1+m_2)/R^3$, $G(m_1+m_2)=\dfrac{\text{DU}^3}{\text{TU}^2}$. Therefore
\begin{equation}
\boxed{\;x_1=-\mu\ \text{DU},\quad x_2=(1-\mu)\ \text{DU};\qquad
Gm_1=(1-\mu)\,\frac{\text{DU}^3}{\text{TU}^2},\quad Gm_2=\mu\,\frac{\text{DU}^3}{\text{TU}^2}\;}
\end{equation}
and the distance between the primaries is $x_2-x_1=1$ DU. This is why the equations of motion contain $(1-\mu)$ and $\mu$ as
gravitational parameters and $(x+\mu)$, $(x+\mu-1)$ in the distances. Since $\mu<\frac12$ the origin is closer to $m_1$.
Earth--Moon: $\mu=1/82.27=0.01216$, so the barycentre is at $\mu R\approx4670$\,km from the Earth's centre, i.e.\ inside the Earth
\notinnotes.

\begin{center}\includegraphics[width=0.36\linewidth]{synodic}\end{center}

\begin{keyf}
$\omega=\sqrt{G(m_1+m_2)/R^3}$; \ $x_1=-\dfrac{m_2}{M}R=-\mu$, \ $x_2=\dfrac{m_1}{M}R=1-\mu$; \ $\mu=\dfrac{m_2}{m_1+m_2}$; \
$Gm_1=1-\mu$, $Gm_2=\mu$ (DU$^3$/TU$^2$).
\end{keyf}

\pitfalls
\begin{itemize}
\item The origin is the centre of mass, \emph{not} the centre of $m_1$.
\item $\mu$ is the mass fraction of the \emph{smaller} primary; $m_1$ is at $-\mu$, $m_2$ at $1-\mu$.
\item Show the two equations (distance and centre of mass) and the step $G(m_1+m_2)=1$ in canonical units.
\end{itemize}

\source{cap05.tex l.111--236 (problem of two bodies, $\omega^2=GM/(a_1+a_2)^3$), l.244--305 (CR3BP, synodic frame, positions and
gravitational parameters).}

% ---------------------------------------------------------------------
\question{B6}{Integrals of motion of the $N$-body problem; the Laplace plane}
% source: cap05.tex l.28-109
\begin{qbox}
Identify the integrals of motion in the problem of $N$ bodies ($\underline R_i$ and $\dot{\underline R}_i$ are the position and velocity of
each point mass $m_i$ in a generic inertial frame) and provide their expressions. Prove that the angular momentum is constant.
Define the invariant Laplace plane and provide the expression of the remarkable vector orthogonal to this plane.
\qmeta{$\times$2 in the list (2025); Jun 2025; Sep 2025 (only: Laplace plane and the vector orthogonal to it).}
\qnote{If only the Laplace plane is asked (Sep 2025): write the equations of motion, the definition of $\underline H_c$, the proof
that it is constant and the definition of the plane (points 1, 4, 5 of the outline).}
\end{qbox}

\begin{outline}
\begin{enumerate}
\item Equations: $m_i\ddot{\underline R}_i=\sum_{j\neq i}Gm_im_j\underline r_{ij}/r_{ij}^3$; internal forces sum to zero.
\item Linear momentum $\sum m_i\dot{\underline R}_i=\underline a$ and centre of mass $\sum m_i\underline R_i=\underline at+\underline b$ (6 integrals).
\item Energy $\mathcal E=T+U$ constant (1 integral).
\item Angular momentum about $C$: $\underline H_c=\sum\underline r_i\times m_i\dot{\underline R}_i$; proof $\dot{\underline H}_c=\underline 0$ (3 integrals).
\item Laplace (invariant) plane: through $C$, orthogonal to $\underline H_c$. Total: 10 integrals.
\end{enumerate}
\end{outline}

\answer
\paragraph{Equations of motion.} $N$ point masses $m_i$, positions $\underline R_i$ in an inertial frame centred in $O$; only mutual
gravitational attraction:
\begin{equation}
m_{i}\,\ddot{\underline R}_{i}=G\sum_{j=1,\,j\neq i}^{N}\frac{m_{i}m_{j}}{r_{ij}^{3}}\,\underline{r}_{ij},\qquad
\underline r_{ij}=\underline R_j-\underline R_i,\ \ r_{ij}=|\underline r_{ij}|.
\end{equation}
The gravitational forces are \emph{internal} and act in equal and opposite pairs ($\underline r_{ji}=-\underline r_{ij}$), so their vector
sum is zero.

\paragraph{1. Linear momentum and centre of mass (6 scalar integrals).} Summing over $i$:
\begin{equation}
\sum_{i=1}^{N}m_i\ddot{\underline R}_i=\underline 0\;\Longrightarrow\;
\sum_{i=1}^{N}m_i\dot{\underline R}_i=\underline a,\qquad \sum_{i=1}^{N}m_i\underline R_i=\underline a\,t+\underline b,
\end{equation}
with $\underline a$, $\underline b$ constant vectors fixed by the initial conditions. The centre of mass
$\underline R_c=\frac{1}{M}\sum m_i\underline R_i$ ($M=\sum m_i$) moves in uniform rectilinear motion: $\underline R_c(t)=(\underline at+\underline b)/M$.

\paragraph{2. Energy (1 scalar integral).} Gravity is conservative, so the total energy (kinetic plus potential,
$U=-\sum_{i<j}Gm_im_j/r_{ij}$) is constant:
\begin{equation}
\mathcal E=\frac12\sum_{i=1}^{N}m_i\,\dot{\underline R}_i\cdot\dot{\underline R}_i-\sum_{i<j}G\frac{m_im_j}{r_{ij}}=\text{const}.
\end{equation}

\paragraph{3. Angular momentum about the centre of mass (3 scalar integrals).}
\begin{equation}
\underline H_c=\sum_{i=1}^{N}\underline r_i\times m_i\dot{\underline R}_i,\qquad \underline r_i=\underline R_i-\underline R_c .
\end{equation}
\emph{Proof that it is constant}:
\begin{equation}
\dot{\underline H}_c=\sum_i(\dot{\underline R}_i-\dot{\underline R}_c)\times m_i\dot{\underline R}_i+\sum_i(\underline R_i-\underline R_c)\times m_i\ddot{\underline R}_i .
\end{equation}
\begin{itemize}
\item First sum: $\sum_i\dot{\underline R}_i\times m_i\dot{\underline R}_i=\underline 0$; the rest is $-\dot{\underline R}_c\times\sum_im_i\dot{\underline R}_i=-\dot{\underline R}_c\times M\dot{\underline R}_c=\underline 0$.
\item Second sum: $-\underline R_c\times\sum_im_i\ddot{\underline R}_i=\underline 0$ by point 1; the remaining term is
\begin{equation}
\sum_i\underline R_i\times m_i\ddot{\underline R}_i=\sum_i\sum_{j\neq i}G\frac{m_im_j}{r_{ij}^3}\,\underline R_i\times(\underline R_j-\underline R_i)
=\sum_i\sum_{j\neq i}G\frac{m_im_j}{r_{ij}^3}\,\underline R_i\times\underline R_j=\underline 0,
\end{equation}
because the coefficient is symmetric in $(i,j)$ while $\underline R_i\times\underline R_j=-\underline R_j\times\underline R_i$: the terms cancel in pairs
(the mutual forces are central, along the line joining the two masses).
\end{itemize}
Hence $\dot{\underline H}_c=\underline 0$: $\underline H_c$ is conserved.

\paragraph{Laplace (invariant) plane.} It is the plane through the centre of mass $C$ \textbf{orthogonal to $\underline H_c$}. Its normal
\begin{equation}
\hat n=\frac{\underline H_c}{|\underline H_c|},\qquad
\underline H_c=\sum_{i=1}^{N}(\underline R_i-\underline R_c)\times m_i\dot{\underline R}_i
\end{equation}
is fixed in inertial space whatever the subsequent motion of the bodies: the plane is \emph{invariable}. It is the natural
reference plane of a planetary system (for the Solar System it is close to the orbit plane of Jupiter \notinnotes).
\begin{center}\includegraphics[width=0.42\linewidth]{nbody}\end{center}

\paragraph{Count.} $\underline a$ (3) $+$ $\underline b$ (3) $+$ $\mathcal E$ (1) $+$ $\underline H_c$ (3) $=$ \textbf{10 scalar integrals}, fixed
by the initial conditions and preserved regardless of the evolution of each individual mass. They are not enough to
integrate the problem for $N\ge3$ (the motion is in general chaotic).

\begin{keyf}
$\sum m_i\dot{\underline R}_i=\underline a$, \ $\sum m_i\underline R_i=\underline at+\underline b$, \ $\mathcal E=\frac12\sum m_i\dot R_i^2-\sum_{i<j}\frac{Gm_im_j}{r_{ij}}$,
\ $\underline H_c=\sum(\underline R_i-\underline R_c)\times m_i\dot{\underline R}_i$; Laplace plane $\perp\underline H_c$ through $C$; 10 integrals.
\end{keyf}

\pitfalls
\begin{itemize}
\item In the proof, show \emph{why} the double sum vanishes (symmetric coefficients, antisymmetric cross product).
\item The Laplace plane is orthogonal to $\underline H_c$ and passes through the centre of mass.
\item Do not confuse $\underline r_i=\underline R_i-\underline R_c$ (position w.r.t.\ $C$) with $\underline r_{ij}$ (mutual position).
\end{itemize}

\further
The angular momentum about the inertial origin, $\underline H_O=\underline H_c+M\underline R_c\times\dot{\underline R}_c$, is also constant, since
$\underline R_c\times\dot{\underline R}_c=(\underline b\times\underline a)/M^2$ \notinnotes. For $N=2$ these integrals reduce the problem to a
Keplerian relative motion (\qref{E10}).

\source{cap05.tex l.28--109 (problem of $N$ bodies: integrals, proof, Laplace plane).}
